%% ══════════════════════════════════════════════════════════════
%%  CHAPTER 8 — PAPER GENERATION ENGINE
%% ══════════════════════════════════════════════════════════════
\chapter{Paper Generation Engine}
\label{ch:paper-engine}

\section{Overview}
\label{sec:paper:overview}

The paper generation engine transforms research analysis into compiled \LaTeX{} PDFs.
It encompasses template selection, section writing, content sanitization, multi-pass compilation, iterative error repair, and reproducible project packaging.
The engine is distributed across 6 modules in \file{server/services/latex/} totaling approximately 2,870 lines.

\section{Text Generation}
\label{sec:paper:text}

\subsection{Section-by-Section Writing}

Papers are written one section at a time via the \texttt{write\_paper\_section} tool.
Claude receives detailed instructions for each section type, including minimum word counts, required subsections, and specific content expectations.

\subsubsection{Empirical Paper Sections}

\begin{description}
  \item[Introduction] (untitled per JF convention): Must present the research question, motivate its importance, preview the methodology, summarize key findings, and outline the paper's contribution relative to existing literature. Target: 2,000--3,000 words.

  \item[Literature Review:] Organized thematically (not chronologically). Must cover 15--30 papers, identify research gaps, and position the paper. Target: 3,000--5,000 words.

  \item[Methodology:] Formal model specification with numbered equations, estimation strategy, identification approach, and assumptions. Must include economic intuition for each variable. Target: 3,000--5,000 words.

  \item[Data:] Source description, sample construction, variable definitions, summary statistics table, preliminary analysis (correlation matrix, unit root tests). Target: 2,000--4,000 words.

  \item[Results:] Main regression tables, robustness checks (alternative specifications, subsamples, placebo tests), figure interpretation. Target: 4,000--7,000 words.

  \item[Discussion:] Economic interpretation, comparison with existing literature, implications for practitioners and policymakers. Target: 2,000--3,000 words.

  \item[Conclusion:] Summary of contributions, limitations, and directions for future research. Target: 1,000--2,000 words.
\end{description}

\subsubsection{Section Rejection and Retry}

The tool executor tracks section rejection attempts.
If a section is rejected more than 3 times (typically due to insufficient length or quality), the system accepts the best available version to prevent infinite loops.

\subsection{Citation Handling}

Citations are managed through BibTeX:

\begin{enumerate}
  \item During the \gauss phase, academic search results include auto-generated BibTeX entries.
  \item Claude accumulates citations throughout the paper, using \texttt{\textbackslash cite\{key\}} and \texttt{\textbackslash citep\{key\}} commands.
  \item The BibTeX file is assembled from all accumulated entries.
  \item \texttt{natbib} with author-year style handles citation formatting.
  \item Duplicate cite keys are resolved by appending letters (e.g., \texttt{smith2020a}, \texttt{smith2020b}).
\end{enumerate}

\subsection{Figure Integration}

Figures are referenced in \LaTeX{} via:

\begin{lstlisting}[language=TeX, caption=Typical figure inclusion pattern.]
\begin{figure}[H]
  \centering
  \includegraphics[width=0.9\textwidth]{fig_01.png}
  \caption{Conditional volatility estimates from GARCH(1,1) model.}
  \label{fig:garch-vol}
\end{figure}
\end{lstlisting}

Figure files are extracted from base64-encoded Python execution results and written to the compilation directory.

\section{LaTeX Templates}
\label{sec:paper:templates}

\subsection{Common Preamble}

The shared preamble (\func{autoScholarPreamble}, $\sim$200 lines) includes 50+ packages:

\begin{multicols}{3}
\begin{itemize}
  \item \pkg{amsmath}
  \item \pkg{amssymb}
  \item \pkg{amsthm}
  \item \pkg{mathtools}
  \item \pkg{geometry}
  \item \pkg{graphicx}
  \item \pkg{booktabs}
  \item \pkg{longtable}
  \item \pkg{multirow}
  \item \pkg{float}
  \item \pkg{natbib}
  \item \pkg{hyperref}
  \item \pkg{cleveref}
  \item \pkg{caption}
  \item \pkg{subcaption}
  \item \pkg{enumitem}
  \item \pkg{microtype}
  \item \pkg{setspace}
  \item \pkg{algorithm2e}
  \item \pkg{listings}
  \item \pkg{xcolor}
  \item \pkg{adjustbox}
  \item \pkg{rotating}
  \item \pkg{dcolumn}
  \item \pkg{threeparttable}
\end{itemize}
\end{multicols}

\subsection{Article Template}

The default template (\func{buildArticleTemplate}):
\begin{itemize}
  \item 12pt, A4, one-column.
  \item 1-inch margins with 1.5 line spacing.
  \item Running headers with abbreviated title.
  \item Numbered sections, Roman numeral tables (per JF convention).
  \item Author block: ``Euler --- The AutoScholar Research Agent.''
  \item Agent disclosure section in footnotes.
\end{itemize}

\subsection{Working Paper Template}

\func{buildWorkingPaperTemplate}:
\begin{itemize}
  \item Single-column with wider margins.
  \item Working paper header.
  \item Draft watermark option.
  \item Otherwise identical preamble.
\end{itemize}

\subsection{Econometrica Template}

\func{buildEconometricaTemplate}:
\begin{itemize}
  \item Two-column layout mimicking \emph{Econometrica} style.
  \item Serif fonts, narrow margins.
  \item Custom theorem environments (Assumption, Proposition, Theorem, Corollary).
\end{itemize}

\section{Content Sanitization}
\label{sec:paper:sanitization}

The sanitizer (\file{latex/sanitizer.ts}, 670 lines) addresses common issues in LLM-generated \LaTeX{}:

\subsection{Special Character Escaping}

\begin{table}[H]
\centering
\caption{\LaTeX{} special character handling.}
\label{tab:latex-escape}
\begin{tabular}{ll}
\toprule
\textbf{Character} & \textbf{Treatment} \\
\midrule
\texttt{\%} & Escaped to \texttt{\textbackslash\%} outside comments \\
\texttt{\&} & Escaped to \texttt{\textbackslash\&} outside tabular environments \\
\texttt{\$} & Context-aware (preserved in math mode, escaped outside) \\
\texttt{\#} & Escaped to \texttt{\textbackslash\#} \\
\texttt{\_} & Escaped to \texttt{\textbackslash\_} outside math mode \\
Greek letters & Wrapped in \texttt{\$...\$} if found outside math mode \\
Em/en dashes & Normalized to \texttt{---}/\texttt{--} \\
\bottomrule
\end{tabular}
\end{table}

\subsection{Table Auto-Fitting}

Wide tables are automatically handled by three progressive strategies:

\begin{enumerate}
  \item \textbf{Column counting:} Tables with $>$6 columns are flagged.
  \item \textbf{\texttt{adjustbox} wrapping:} First attempt wraps the table in \texttt{\textbackslash begin\{adjustbox\}\{max width=\textbackslash textwidth\}}.
  \item \textbf{\texttt{resizebox} fallback:} If \texttt{adjustbox} fails, uses \texttt{\textbackslash resizebox\{\textbackslash textwidth\}\{!\}}.
  \item \textbf{Safety wrapping:} Final pass adds overflow protection to all remaining tables.
\end{enumerate}

\subsection{Blank Figure Detection}

\func{isBlankPng} detects empty figures at the binary level:

\begin{enumerate}
  \item Reads the PNG file and locates IDAT (image data) chunks.
  \item Decompresses the raw pixel data via zlib inflate.
  \item Strips PNG filter bytes (one per row).
  \item Computes pixel variance across the image.
  \item If variance $< 50.0$, the figure is classified as blank and excluded.
\end{enumerate}

\section{Compilation Pipeline}
\label{sec:paper:compilation}

\subsection{Binary Discovery}

The system searches for \texttt{pdflatex} and \texttt{bibtex} across multiple macOS installation paths:

\begin{itemize}
  \item \texttt{/Library/TeX/texbin/}
  \item \texttt{/usr/local/texlive/2024/bin/universal-darwin/}
  \item \texttt{/usr/texbin/}
  \item System PATH
\end{itemize}

\subsection{Multi-Pass Compilation}

Standard academic compilation requires multiple passes:

\begin{enumerate}
  \item \textbf{Pass 1:} \texttt{pdflatex} --- Initial compilation, generates \texttt{.aux} files.
  \item \textbf{Pass 2:} \texttt{bibtex} --- Processes \texttt{.aux} to resolve citations from \texttt{.bib}.
  \item \textbf{Pass 3:} \texttt{pdflatex} --- Re-processes with resolved citations.
  \item \textbf{Pass 4:} \texttt{pdflatex} --- Final pass to resolve all cross-references.
\end{enumerate}

All passes run with \texttt{-interaction=nonstopmode} to prevent interactive prompts.

\subsection{Error Classification}

The \LaTeX{} log parser (\file{latexLogParser.ts}, 146 lines) classifies errors into 8 types:

\begin{description}
  \item[Undefined control sequence:] Misspelled commands or missing packages.
  \item[Math errors:] Mismatched delimiters, missing \$ signs.
  \item[Unicode errors:] Non-ASCII characters in text mode.
  \item[Environment errors:] Mismatched \texttt{\textbackslash begin}/\texttt{\textbackslash end} pairs.
  \item[Package errors:] Package conflicts or missing options.
  \item[Fatal errors:] Emergency stops, capacity exceeded.
  \item[Missing characters:] Font encoding issues.
  \item[Other:] Unclassified errors.
\end{description}

Additionally tracks overfull/underfull box warnings (hbox/vbox).

\subsection{Iterative Repair Loop}

When compilation fails, the \LaTeX{} Patch Manager (\file{latexPatchManager.ts}, 369 lines) orchestrates repair:

\begin{algorithm}[H]
\caption{\LaTeX{} Iterative Repair}
\label{alg:latex-repair}
\SetAlgoLined
\KwIn{Failed compilation, max attempts = 8}
\KwOut{Compiled PDF or final failure}

Register failed compilation metadata\;
\For{$i \gets 1$ \KwTo 8}{
  $\text{errors} \gets \func{inspectLatexErrors}()$\;
  Format errors for Claude with line numbers and context\;
  Claude generates a diff patch\;
  $\text{result} \gets \func{patchLatex}(\text{diff})$\;
  \If{patch failed}{
    Try search-replace fallback\;
    Try trimmed matching fallback\;
  }
  $\text{pdf} \gets \func{recompileLatex}()$\;
  \If{compilation succeeded}{
    \Return{pdf}\;
  }
}
\Return{failure with last error log}\;
\end{algorithm}

\section{Document Export}
\label{sec:paper:export}

Beyond \LaTeX{} PDF, the system supports three additional export formats:

\subsection{PDF via Puppeteer}

\file{pdf/puppeteerPdf.ts} (75 lines):
\begin{enumerate}
  \item Converts markdown to HTML with KaTeX math rendering.
  \item Loads in headless Chromium.
  \item Waits for math rendering completion.
  \item Generates A4 PDF with AutoScholar-branded footer.
\end{enumerate}

\subsection{DOCX via Pandoc}

\file{pdf/pandocDocx.ts} (50 lines):
\begin{enumerate}
  \item Writes markdown to temp file.
  \item Invokes \texttt{pandoc -f markdown -t docx}.
  \item Returns DOCX buffer.
\end{enumerate}

\subsection{HTML Export}

\file{pdf/htmlExport.ts} (9 lines):
\begin{itemize}
  \item Self-contained HTML with embedded CSS, KaTeX, and syntax highlighting.
  \item Responsive design with print-optimized CSS.
  \item AutoScholar branding with embedded SVG logo.
\end{itemize}

\section{Reproducible Project Packaging}
\label{sec:paper:packaging}

The packaging module (\file{latex/packaging.ts}, 358 lines) creates ZIP archives following the \gls{rre} structure:

\subsection{Directory Structure}

\begin{verbatim}
project/
  paper.tex              Main LaTeX source
  references.bib         BibTeX references
  outputs/
    paper.pdf            Compiled PDF
    figures/             All generated figures (PNG)
  code/
    01_data_download.py  Pipeline-ordered scripts
    02_data_prep.py
    03_analysis.py
    ...
  data/
    raw/api/             API-sourced data (CSV, JSON)
    raw/web/             Web-extracted data
    processed/           Cleaned datasets
    metadata/            Data provenance (JSON sidecars)
  run_pipeline.py        Auto-generated master script
  requirements.txt       Auto-detected dependencies
  README.md              Build instructions
\end{verbatim}

\subsection{Automatic Dependency Detection}

The \texttt{requirements.txt} generator:
\begin{enumerate}
  \item Scans all Python scripts for \texttt{import} and \texttt{from ... import} statements.
  \item Maps module names to pip package names (e.g., \texttt{sklearn} $\to$ \texttt{scikit-learn}, \texttt{cv2} $\to$ \texttt{opencv-python}).
  \item Deduplicates and sorts alphabetically.
  \item Excludes standard library modules.
\end{enumerate}

\subsection{Master Pipeline Script}

\texttt{run\_pipeline.py} is auto-generated with:
\begin{enumerate}
  \item Sequential execution of all \texttt{code/*.py} scripts in numeric order.
  \item Error handling with script-level try/catch.
  \item Summary output listing successful/failed scripts.
  \item Working directory setup for data and output paths.
\end{enumerate}
